\ifdefined\gkver\else\def\gkver{student}\fi
\documentclass[\gkver]{gaokaozhenti}
\begin{document}

\chapter{2021年江苏}

\begin{enumerate}
\item
%% number: 1
%% paperName: 2021年江苏高考第 1 题 4 分
%% typeId: 1
%% type: 单选题
%% chapter: 原子和原子核
%% point: 人工核反应
%% method: 无
%% score: 4
%% degree: 950
%% duplicateId: 0
%% body:
用“中子活化”技术分析某样品的成分，中子轰击样品中的 \ce{^{14}_{7}N} 产生 \ce{^{14}_{6}C} 和另一种粒子 X，则 X 是\xzanswer{A}

\fourchoices[answer=A]
{质子}
{$\alpha$粒子}
{$\beta$粒子}
{正电子}

%% answer: A
%% memo:
\memoanswer{
该核反应方程为 \ce{^{14}_{7}N + ^{1}_{0}n -> ^{14}_{6}C + ^{1}_{1}H}，可知 X 是质子。

故选 A。
}

\item
%% number: 2
%% paperName: 2021年江苏高考第 2 题 4 分
%% typeId: 1
%% type: 单选题
%% chapter: 电路
%% point: 电容
%% method: 无
%% score: 4
%% degree: 850
%% duplicateId: 0
%% body:
有研究发现，某神经细胞传递信号时，离子从细胞膜一侧流到另一侧形成跨膜电流，若将该细胞膜视为 $1\times10^{-8}\UF$ 的电容器，在 $2\UmsT$ 内细胞膜两侧的电势差从 $-70\UmV$ 变为 $30\UmV$，则该过程中跨膜电流的平均值为\xzanswer{D}

\fourchoices[answer=D]
{$1.5\times10^{-7}\UA$}
{$2\times10^{-7}\UA$}
{$3.5\times10^{-7}\UA$}
{$5\times10^{-7}\UA$}

%% answer: D
%% memo:
\memoanswer{
根据 $Q = CU$ 可知：$\Delta Q = C\Delta U = 10^{-8}\UF \times (30 + 70)\UmV = 10^{-9}\UC$。

则该过程中跨膜电流的平均值为 $I = \frac{\Delta Q}{t} = \frac{10^{-9}}{2\times10^{-3}}\UA = 5\times10^{-7}\UA$。

故选 D。
}

\item
%% number: 3
%% paperName: 2021年江苏高考第 3 题 4 分
%% typeId: 1
%% type: 单选题
%% chapter: 曲线运动
%% point: 人造地球卫星
%% method: 无
%% score: 4
%% degree: 850
%% duplicateId: 0
%% body:
我国航天人发扬“两弹一星”精神砥砺前行，从“东方红一号”到“北斗”不断创造奇迹。“北斗”第 49 颗卫星的发射迈出组网的关键一步。该卫星绕地球做圆周运动，运动周期与地球自转周期相同，轨道平面与地球赤道平面成一定夹角。该卫星\xzanswer{B}

\fourchoices[answer=B]
{运动速度大于第一宇宙速度}
{运动速度小于第一宇宙速度}
{轨道半径大于“静止”在赤道上空的同步卫星}
{轨道半径小于“静止”在赤道上空的同步卫星}

%% answer: B
%% memo:
\memoanswer{
AB．第一宇宙速度是指绕地球表面做圆周运动的速度，是环绕地球做圆周运动的所有卫星的最大环绕速度，该卫星的运转半径远大于地球的半径，可知运行线速度小于第一宇宙速度，选项 A 错误 B 正确；

CD．根据 $G\frac{Mm}{r^{2}}=m\frac{4\pi^{2}}{T^{2}}r$ 可知 $r=\sqrt[3]{\frac{GMT^{2}}{4\pi^{2}}}$，因为该卫星的运动周期与地球自转周期相同，等于“静止”在赤道上空的同步卫星的周期，可知该卫星的轨道半径等于“静止”在赤道上空的同步卫星的轨道半径，选项 CD 错误。

故选 B。
}

\item
%% number: 4
%% paperName: 2021年江苏高考第 4 题 4 分
%% typeId: 1
%% type: 单选题
%% chapter: 机械振动
%% point: 简谐振动
%% method: 无
%% score: 4
%% degree: 650
%% duplicateId: 0
%% body:
如图所示，半径为 $R$ 的圆盘边缘有一钉子 B，在水平光线下，圆盘的转轴 A 和钉子 B 在右侧墙壁上形成影子 O 和 P，以 O 为原点在竖直方向上建立 $x$ 坐标系。$t=0$ 时从图示位置沿逆时针方向匀速转动圆盘，角速度为 $\omega$，则 P 做简谐运动的表达式为\xzanswer{B}

\onepicture[w=6.5cm]{figs/04.png}

\fourchoices[answer=B]
{$x = R\sin(\omega t - \frac{\pi}{2})$}
{$x = R\sin(\omega t + \frac{\pi}{2})$}
{$x = 2R\sin(\omega t - \frac{\pi}{2})$}
{$x = 2R\sin(\omega t + \frac{\pi}{2})$}

%% answer: B
%% memo:
\memoanswer{
由图可知，影子 P 做简谐运动的振幅为 $R$，以向上为正方向，设 P 的振动方程为 $x = R\sin(\omega t + \varphi)$。

由图可知，当 $t=0$ 时，P 的位移为 $R$，所用时间为 $t=\frac{\pi}{\omega}$。

代入振动方程解得 $\varphi=\frac{\pi}{2}$。

则 P 做简谐运动的表达式为 $x = R\sin(\omega t + \frac{\pi}{2})$。

故 B 正确，ACD 错误。

故选 B。
}

\item
%% number: 5
%% paperName: 2021年江苏高考第 5 题 4 分
%% typeId: 1
%% type: 单选题
%% chapter: 磁场
%% point: 安培力
%% method: 等效替代
%% score: 4
%% degree: 750
%% duplicateId: 0
%% body:
在光滑桌面上将长为 $\pi L$ 的软导线两端固定，固定点的距离为 $2L$，导线通有电流 $I$，处于磁感应强度大小为 $B$、方向竖直向下的匀强磁场中，导线中的张力为\xzanswer{A}

\fourchoices[answer=A]
{$BIL$}
{$2BIL$}
{$\pi BIL$}
{$2\pi BIL$}

%% answer: A
%% memo:
\memoanswer{
从上向下看导线的图形如图所示。

\onepicture[w=3.2cm]{figs/05_an.png}

导线的有效长度为 $2L$，则所受的安培力大小 $F = 2BIL$。

设绳子的张力为 $T$，由几何关系可知 $T = \frac{F}{2}$。

解得 $T = BIL$。

故 A 正确，BCD 错误。

故选 A。
}

\item
%% number: 6
%% paperName: 2021年江苏高考第 6 题 4 分
%% typeId: 1
%% type: 单选题
%% chapter: 光学
%% point: 光的干涉
%% method: 无
%% score: 4
%% degree: 750
%% duplicateId: 0
%% body:
铁丝圈上附有肥皂膜，竖直放置时，肥皂膜上的彩色条纹上疏下密，由此推测肥皂膜前后两个面的侧视形状应当是\xzanswer{C}

\fourchoices[answer=C, ispicture=true, h=2.2cm]
{figs/06a.png}
{figs/06b.png}
{figs/06c.png}
{figs/06d.png}

%% answer: C
%% memo:
\memoanswer{
薄膜干涉为前后两个面反射回来的光发生干涉形成干涉条纹，在复色光时，出现彩色条纹，由于重力作用，肥皂膜前后表面的厚度从上到下逐渐增大，从而使干涉条纹的间距上密下疏，由于表面张力的作用，使得肥皂膜向内凹陷，故 C 正确，ABD 错误。

故选 C。
}

\item
%% number: 7
%% paperName: 2021年江苏高考第 7 题 4 分
%% typeId: 1
%% type: 单选题
%% chapter: 光学
%% point: 全反射
%% method: 无
%% score: 4
%% degree: 750
%% duplicateId: 0
%% body:
某种材料制成的半圆形透明砖平放在方格纸上，将激光束垂直于 AC 面射入，可以看到光束从圆弧面 ABC 出射，沿 AC 方向缓慢平移该砖，在如图所示位置时，出射光束恰好消失，该材料的折射率为\xzanswer{A}

\onepicture[w=4.5cm]{figs/07.png}

\fourchoices[answer=A]
{1.2}
{1.4}
{1.6}
{1.8}

%% answer: A
%% memo:
\memoanswer{
画出激光束从玻璃砖射出时恰好发生全反射的入射角如图所示。

\onepicture[w=5.0cm]{figs/07_an.png}

全反射的条件：$\sin\theta=\frac{1}{n}$。

由几何关系知：$\sin\theta=\frac{5}{6}$。

联立解得：$n=1.2$。

故 A 正确，BCD 错误。

故选 A。
}

\item
%% number: 8
%% paperName: 2021年江苏高考第 8 题 4 分
%% typeId: 1
%% type: 单选题
%% chapter: 光学
%% point: 光电效应
%% method: 无
%% score: 4
%% degree: 650
%% duplicateId: 0
%% body:
如图所示，分别用 1、2 两种材料作 K 极进行光电效应探究，其截止频率 $\nu_{1}<\nu_{2}$，保持入射光不变，则光电子到达 A 极时动能的最大值 $E_{km}$ 随电压 $U$ 变化关系的图像是\xzanswer{C}

\onepicture[w=4.0cm]{figs/08a.png}

\fourchoices[answer=C, ispicture=true, h=2.2cm]
{figs/08b.png}
{figs/08c.png}
{figs/08d.png}
{figs/08e.png}

%% answer: C
%% memo:
\memoanswer{
光电管所加电压为正向电压，则根据爱因斯坦光电效应方程可知光电子到达 A 极时动能的最大值

$E_{km}=Ue+h\nu-h\nu_{\text{截止}}$

可知 $E_{km}-U$ 图像的斜率相同，均为 $e$；截止频率越大，则图像在纵轴上的截距越小，因 $\nu_{1}<\nu_{2}$，则图像 C 正确，ABD 错误。

故选 C。
}

\item
%% number: 9
%% paperName: 2021年江苏高考第 9 题 4 分
%% typeId: 1
%% type: 单选题
%% chapter: 曲线运动
%% point: 斜抛运动
%% method: 无
%% score: 4
%% degree: 550
%% duplicateId: 0
%% body:
如图所示，A、B 两篮球从相同高度同时抛出后直接落入篮筐，落入篮筐时的速度方向相同，下列判断正确的是\xzanswer{D}

\onepicture[w=6.0cm]{figs/09.png}

\fourchoices[answer=D]
{A 比 B 先落入篮筐}
{A、B 运动的最大高度相同}
{A 在最高点的速度比 B 在最高点的速度小}
{A、B 上升到某一相同高度时的速度方向相同}

%% answer: D
%% memo:
\memoanswer{
AB．若研究两个过程的逆过程，可看做是从篮筐沿同方向斜向上的斜抛运动，落到同一高度上的 A、B 两点，则 A 上升的高度较大，高度决定时间，可知 A 运动时间较长，即 B 先落入篮筐中，AB 错误；

C．因为两球抛射角相同，A 的射程较远，则 A 球的水平速度较大，即在最高点的速度比 B 在最高点的速度大，C 错误；

D．由斜抛运动的对称性可知，当 A、B 上升到与篮筐相同高度时的速度方向相同，D 正确。

故选 D。
}

\item
%% number: 10
%% paperName: 2021年江苏高考第 10 题 4 分
%% typeId: 1
%% type: 单选题
%% chapter: 电场
%% point: 电场强度
%% method: 对称法
%% score: 4
%% degree: 450
%% duplicateId: 0
%% body:
一球面均匀带有正电荷，球内的电场强度处处为零，如图所示，O 为球心，A、B 为直径上的两点，$OA = OB$，现垂直于 AB 将球面均分为左右两部分，C 为截面上的一点，移去左半球面，右半球面所带电荷仍均匀分布，则\xzanswer{A}

\onepicture[w=4.5cm]{figs/10.png}

\fourchoices[answer=A]
{O、C 两点电势相等}
{A 点的电场强度大于 B 点}
{沿直线从 A 到 B 电势先升高后降低}
{沿直线从 A 到 B 电场强度逐渐增大}

%% answer: A
%% memo:
\memoanswer{
A．由于球壳内部的场强为零，补全以后可知在右侧球壳任取一点，该点和 C 连线后过左半边球壳一点，两点对 C 的库仑力为零，即等大反向共线，将左侧沿竖直对称到右侧球面，则可以知道这两个点在 C 点的合场强水平向左，根据叠加原理可知 C 点场强向左，同理 OC 上都是水平向左，因此 OC 是等势线，故 A 正确；

BD．将题中半球壳补成一个完整的球壳，且带电均匀，设左、右半球在 A 点产生的电场强度大小分别为 $E_{1}$ 和 $E_{2}$；由题知，均匀带电球壳内部电场强度处处为零，则知 $E_{1} = E_{2}$，根据对称性，左右半球在 B 点产生的电场强度大小分别为 $E_{2}$ 和 $E_{1}$，且 $E_{1} = E_{2}$。

在图示电场中，A 的电场强度大小为 $E_{2}$，方向向左，B 的电场强度大小为 $E_{1}$，方向向左，所以 A 点的电场强度与 B 点的电场强度相同，沿直线从 A 到 B 电场强度不可能逐渐增大，故 BD 错误；

C．根据电场的叠加原理可知，在 AB 连线上电场线方向向左，沿着电场线方向电势逐渐降低，则沿直线从 A 到 B 电势升高，故 C 错误。

故选 A。
}

\item
%% number: 11
%% paperName: 2021年江苏高考第 11 题 15 分
%% typeId: 4
%% type: 实验题
%% chapter: 运动和力
%% point: 动量守恒定律
%% method: 无
%% score: 15
%% degree: 450
%% duplicateId: 0
%% body:
小明利用如图 \ref{2021江苏11a} 所示的实验装置验证动量定理。将遮光条安装在滑块上，用天平测出遮光条和滑块的总质量 $M = 200.0\Ug$，槽码和挂钩的总质量 $m = 50.0\Ug$。实验时，将滑块系在绕过定滑轮悬挂有槽码的细线上。滑块由静止释放，数字计时器记录下遮光条通过光电门 1 和 2 的遮光时间 $\Delta t_{1}$ 和 $\Delta t_{2}$，以及这两次开始遮光的时间间隔 $\Delta t$，用游标卡尺测出遮光条宽度，计算出滑块经过两光电门速度的变化量 $\Delta v$。

\twopicture[labelA=2021江苏11a, labelB=2021江苏11b, numA=1, numB=2, wA=8.4cm, wB=4.2cm]
{figs/11a.png}{figs/11b.png}

\begin{enumerate}
	\item
	游标卡尺测量遮光条宽度如图 \ref{2021江苏11b} 所示，其宽度 $d =$ \tkanswer{10.20} $\Umm$；
	\item
	打开气泵，待气流稳定后调节气垫导轨，直至看到导轨上的滑块能在短时间内保持静止，其目的是 \tkanswer{将气垫导轨调至水平}；
	\item
	多次改变光电门 2 的位置进行测量，得到 $\Delta t$ 和 $\Delta v$ 的数据如下表，请根据表中数据，在方格纸上作出 $\Delta v-\Delta t$ 图线。

	\begin{center}
	\begin{tblr}{|c|c|c|c|c|c|}
	\hline
	$\Delta t/\Us$ & 0.721 & 0.790 & 0.854 & 0.913 & 0.968 \\
	\hline
	$\Delta v/\Umdsn$ & 1.38 & 1.52 & 1.64 & 1.75 & 1.86 \\
	\hline
	\end{tblr}
	\end{center}

	\onepicture[w=7.0cm]{figs/11c.png}

	\item
	查得当地的重力加速度 $g = 9.8\Umsq$，根据动量定理，$\Delta v-\Delta t$ 图线斜率的理论值为 \tkanswer{$1.92\sim1.97$} $\Umsq$；
	\item
	实验结果发现，图线斜率的实验值总小于理论值，产生这一误差的两个可能原因是 \tkanswer{BD}。

	A．选用的槽码质量偏小

	B．细线与气垫导轨不完全平行

	C．每次释放滑块的位置不同

	D．实验中 $\Delta t$ 的测量值偏大
\end{enumerate}

%% answer:
\jdanswer{
\begin{enumerate}
	\item
	$10.20$

	\item
	将气垫导轨调至水平

	\item
	（作图见解析）

	\item
	$1.92\sim1.97$

	\item
	BD
\end{enumerate}
}
%% memo:
\memoanswer{
（1）游标卡尺的读数为 $10\Umm + 4\times0.05\Umm = 10.20\Umm$。

（2）滑块保持稳定，说明气垫导轨水平。

（3）根据表格中数据描点并用直线连接。

\onepicture[w=7.0cm]{figs/11_an.png}

（4）取直线与表格交点的最大距离可知图象的斜率为 $k = \frac{1.9-1.34}{0.99-0.70}\Umsq = 1.93\Umsq$。

（5）根据动量定理 $F\Delta t = M\Delta v$ 变形得 $\frac{\Delta v}{\Delta t} = \frac{F}{M}$。

A．槽码质量偏小，而实际的槽码质量偏大，则合外力 $F$ 偏大，所以图线斜率的实验值偏大，A 错误；

B．细线与气垫导轨不平行，滑块实际所受合外力为 $F$ 的水平分力，所以图线斜率的实验值偏小，B 正确；

C．滑块释放的位置与斜率相关的参量无关，C 错误；

D．$\Delta t$ 偏大，则 $\frac{\Delta v}{\Delta t}$ 偏小，图线斜率偏小，D 正确。

故选 BD。
}

\item
%% number: 12
%% paperName: 2021年江苏高考第 12 题 8 分
%% typeId: 6
%% type: 计算题
%% chapter: 电磁感应
%% point: 切割磁感线电动势
%% method: 无
%% score: 8
%% degree: 650
%% duplicateId: 0
%% body:
贯彻新发展理念，我国风力发电发展迅猛，2020 年我国风力发电量高达 4000 亿千瓦时。某种风力发电机的原理如图所示，发电机的线圈固定，磁体在叶片驱动下绕线圈对称轴转动，已知磁体间的磁场为匀强磁场，磁感应强度的大小为 $0.20\UT$，线圈的匝数为 100、面积为 $0.5\Umq$，电阻为 $0.6\UO$，若磁体转动的角速度为 $90\Urads$，线圈中产生的感应电流为 $50\UA$。求：
\begin{enumerate}
	\item
	线圈中感应电动势的有效值 $E$；
	\item
	线圈的输出功率 $P$。
\end{enumerate}

\onepicture[w=4.8cm, align=right]{figs/12.png}

%% answer:
\jdanswer{
\begin{enumerate}
	\item
	$E = 6.4\times10^{2}\UV$
	\item
	$P = 3.1\times10^{4}\UW$
\end{enumerate}
}
%% memo:
\memoanswer{
（1）电动势的最大值 $E_{m}=NBS\omega$。

有效值 $E=\frac{E_{m}}{\sqrt{2}}$。

解得 $E=\frac{NBS\omega}{\sqrt{2}}$。

代入数据得 $E=6.4\times10^{2}\UV$。

（2）输出电压 $U=E-Ir$。

输出功率 $P=IU$。

解得 $P=I(E-Ir)$。

代入数据得 $P=3.1\times10^{4}\UW$。
}

\item
%% number: 13
%% paperName: 2021年江苏高考第 13 题 8 分
%% typeId: 6
%% type: 计算题
%% chapter: 气体性质
%% point: 热力学第一定律
%% method: 无
%% score: 8
%% degree: 650
%% duplicateId: 0
%% body:
如图所示，一定质量理想气体被活塞封闭在气缸中，活塞的面积为 $S$，与气缸底部相距 $L$，气缸和活塞绝热性能良好，气体的压强、温度与外界大气相同，分别为 $p_{0}$ 和 $t_{0}$。现接通电热丝加热气体，一段时间后断开，活塞缓慢向右移动距离 $L$ 后停止，活塞与气缸间的滑动摩擦为 $f$，最大静摩擦力等于滑动摩擦力，整个过程中气体吸收的热量为 $Q$，求该过程中
\begin{enumerate}
	\item
	内能的增加量 $\Delta U$；
	\item
	最终温度 $T$。
\end{enumerate}

\onepicture[w=6.5cm, align=right]{figs/13.png}

%% answer:
\jdanswer{
\begin{enumerate}
	\item
	$\Delta U = Q - (p_{0}S + f)L$
	\item
	$T = \frac{2(p_{0}S + f)}{p_{0}S}T_{0}$
\end{enumerate}
}
%% memo:
\memoanswer{
（1）活塞移动时受力平衡 $p_{1}S = p_{0}S + f$。

气体对外界做功 $W = p_{1}SL$。

根据热力学第一定律 $\Delta U = Q - W$。

解得 $\Delta U = Q - (p_{0}S + f)L$。

（2）活塞发生移动前，等容过程 $\frac{p_{0}}{t_{0}}=\frac{p_{1}}{t_{1}}$。

活塞向右移动了 $L$，等压过程 $\frac{V_{1}}{t_{1}}=\frac{V_{2}}{T}$。

且 $V_{2}=2V_{1}$。

解得 $T = \frac{2(p_{0}S + f)}{p_{0}S}T_{0}$。
}

\item
%% number: 14
%% paperName: 2021年江苏高考第 14 题 13 分
%% typeId: 6
%% type: 计算题
%% chapter: 曲线运动
%% point: 圆周运动的应用
%% method: 无
%% score: 13
%% degree: 550
%% duplicateId: 0
%% body:
如图所示的离心装置中，光滑水平轻杆固定在竖直转轴的 O 点，小圆环 A 和轻质弹簧套在轻杆上，长为 $2L$ 的细线和弹簧两端分别固定于 O 和 A，质量为 $m$ 的小球 B 固定在细线的中点，装置静止时，细线与竖直方向的夹角为 $37^{\circ}$，现将装置由静止缓慢加速转动，当细线与竖直方向的夹角增大到 $53^{\circ}$ 时，A、B 间细线的拉力恰好减小到零，弹簧弹力与静止时大小相等、方向相反，重力加速度为 $g$，取 $\sin37^{\circ}=0.6$，$\cos37^{\circ}=0.8$，求：
\begin{enumerate}
	\item
	装置静止时，弹簧弹力的大小 $F$；
	\item
	环 A 的质量 $M$；
	\item
	上述过程中装置对 A、B 所做的总功 $W$。
\end{enumerate}

\onepicture[w=6.5cm, align=right]{figs/14.png}

%% answer:
\jdanswer{
\begin{enumerate}
	\item
	$\frac{3mg}{8}$
	\item
	$\frac{9}{64}m$
	\item
	$\frac{31}{30}mgL$
\end{enumerate}
}
%% memo:
\memoanswer{
（1）设 AB、OB 的张力分别为 $F_{1}$、$F_{2}$，A 受力平衡 $F = F_{1}\sin37^{\circ}$。

B 受力平衡 $F_{1}\cos37^{\circ}+F_{2}\cos37^{\circ}=mg$，$F_{1}\sin37^{\circ}=F_{2}\sin37^{\circ}$。

解得 $F=\frac{3mg}{8}$。

（2）设装置转动的角速度为 $\omega$，对 A 有 $F=M\omega^{2}\cdot\frac{8}{5}L$。

对 B 有 $mg\tan53^{\circ} = m\omega^{2}\cdot\frac{4}{5}L$。

解得 $M = \frac{9}{64}m$。

（3）B 上升的高度 $h = \frac{1}{5}L$，A、B 的动能分别为

$E_{\mathrm{kA}} = \frac{1}{2}M\left(\omega\frac{8}{5}L\right)^{2}$，$E_{\mathrm{kB}} = \frac{1}{2}m\left(\omega\frac{4}{5}L\right)^{2}$。

根据能量守恒定律可知 $W = (E_{\mathrm{kA}} - 0) + (E_{\mathrm{kB}} - 0) + mgh$。

解得 $W = \frac{31}{30}mgL$。
}

\item
%% number: 15
%% paperName: 2021年江苏高考第 15 题 16 分
%% typeId: 6
%% type: 计算题
%% chapter: 磁场
%% point: 带电粒子在磁场中的运动
%% method: 无
%% score: 16
%% degree: 350
%% duplicateId: 0
%% body:
如图 \ref{2021江苏15a} 所示，回旋加速器的圆形匀强磁场区域以 O 点为圆心，磁感应强度大小为 $B$，加速电压的大小为 $U$、质量为 $m$、电荷量为 $q$ 的粒子从 O 附近飘入加速电场，多次加速后粒子经过 P 点绕 O 做圆周运动，半径为 $R$，粒子在电场中的加速时间可以忽略。为将粒子引出磁场，在 P 位置安装一个“静电偏转器”，如图 \ref{2021江苏15b} 所示，偏转器的两极板 M 和 N 厚度均匀，构成的圆弧形狭缝圆心为 Q、圆心角为 $\alpha$，当 M、N 间加有电压时，狭缝中产生电场强度大小为 $E$ 的电场，使粒子恰能通过狭缝，粒子在再次被加速前射出磁场，不计 M、N 间的距离。求：
\begin{enumerate}
	\item
	粒子加速到 P 点所需要的时间 $t$；
	\item
	极板 N 的最大厚度 $d_{m}$；
	\item
	磁场区域的最大半径 $R_{m}$。
\end{enumerate}

\twopicture[labelA=2021江苏15a, labelB=2021江苏15b, numA=1, numB=2, wA=5.0cm, wB=5.0cm, align=right]
{figs/15a.png}{figs/15b.png}

%% answer:
\jdanswer{
\begin{enumerate}
	\item
	$\left(\frac{qB^{2}R^{2}}{2mU} - 1\right)\frac{\pi m}{qB}$
	\item
	$2\left(\sqrt{R^{2} - \frac{2mU}{qB^{2}}} - \sqrt{R^{2} - \frac{4mU}{qB^{2}}}\right)$
	\item
	$R + \frac{2mER}{qB^{2}R - mE}\sin\frac{\alpha}{2}$
\end{enumerate}
}
%% memo:
\memoanswer{
（1）设粒子在 P 的速度大小为 $v_{P}$，则根据 $qvB = m\frac{v^{2}}{r}$ 可知半径表达式为 $R=\frac{mv_{P}}{qB}$。

根据动能定理粒子在静电场中加速有 $nqU=\frac{1}{2}mv_{P}^{2}$。

粒子在磁场中运动的周期为 $T=\frac{2\pi m}{qB}$。

粒子运动的总时间为 $t=(n-1)\times\frac{T}{2}$。

解得 $t=\left(\frac{qB^{2}R^{2}}{2mU}-1\right)\frac{\pi m}{qB}$。

（2）由粒子的运动半径 $r = \frac{mv}{qB}$，结合动能表达式 $E_{k} = \frac{1}{2}mv^{2}$ 变形得 $r = \frac{\sqrt{2mE_{k}}}{qB}$。

则粒子加速到 P 前最后两个半周的运动半径为

$r_{1}=\frac{\sqrt{2m(E_{kP}-qU)}}{qB}$，$r_{2}=\frac{\sqrt{2m(E_{kP}-2qU)}}{qB}$。

由几何关系 $d_{m}=2(r_{1}-r_{2})$。

结合 $E_{kP}=\frac{(qBR)^{2}}{2m}$ 解得

$d_{m}=2\left(\sqrt{R^{2}-\frac{2mU}{qB^{2}}}-\sqrt{R^{2}-\frac{4mU}{qB^{2}}}\right)$。

（3）设粒子在偏转器中的运动半径为 $r_{Q}$，则在偏转器中，要使粒子半径变大，电场力应和洛伦兹力反向，共同提供向心力

$qv_{P}B - qE = m\frac{v_{P}^{2}}{r_{Q}}$。

设粒子离开偏转器的点为 E，圆周运动的圆心为 $O'$。由题意知，$O'$ 在 SQ 上，且粒子飞离磁场的点与 O、$O'$ 在一条直线上。

\onepicture[w=7.5cm]{figs/15_an1.png}

粒子在偏转器中运动的圆心在 Q 点，从偏转器飞出，即从 E 点离开，又进入回旋加速器中的磁场，此时粒子的运动半径又变为 $R$，然后轨迹发生偏离，从偏转器的 F 点飞出磁场，那么磁场的最大半径即为

$R_{\mathrm{m}}=OF=R+OO'$。

将等腰三角形 $\triangle OO'Q$ 放大如图。

\onepicture[w=5.0cm]{figs/15_an2.png}

虚线为从 Q 点向 $OO'$ 所引垂线，虚线平分 $\alpha$ 角，则

$OO'=2(r_{Q}-R)$。

解得最大半径为

$R_{\mathrm{m}}=R+\frac{2mER}{qB^{2}R-mE}\sin\frac{\alpha}{2}$。
}

\end{enumerate}

\end{document}
